\subsection{最初に必要な用語}

\par\smallskip\noindent\refstepcounter{table}\label{tab:ch16-01}表 \thetable: 最初に必要な用語における用語・初学者向けの説明\par\smallskip
表 \thetable は、最初に必要な用語における用語・初学者向けの説明を比較・確認しやすい形で整理したものである。\par\smallskip
\begin{longtable}{>{\raggedright\arraybackslash}p{0.26\textwidth}>{\raggedright\arraybackslash}p{0.68\textwidth}}
\toprule
用語 & 初学者向けの説明 \\
\midrule
\endfirsthead
\toprule
用語 & 初学者向けの説明 \\
\midrule
\endhead
\term{コンピューティング}{Computing} & 計算機を使って情報を表し、処理し、保存し、通信する活動全体である。 \\
システム & 目的を達成するために、ソフトウェア、ハードウェア、人、手順、データなどが組み合わさったものである。 \\
\term{サブシステム}{Subsystem} & システムを分けた部分であり、独立した役割を持つ。さらにモジュールや部品へ分けられる。 \\
\term{モジュール}{Module} & 一つの責務を持つように切り分けたソフトウェアの単位である。変更や理解の単位になりやすい。 \\
\term{コンピュータアーキテクチャ}{Computer Architecture} & コンピュータが外から見て何をするか、どのような構成要素を持つかを表す考え方である。 \\
\term{コンピュータ組織}{Computer Organization} & 構成要素がどのようにつながり、命令をどのように実行するかという内部の実装構成である。 \\
\term{データ構造}{Data Structure} & データを効率よく扱うための並べ方、結び方、格納の仕方である。 \\
\term{アルゴリズム}{Algorithm} & 問題を解くための手順であり、入力を受け取り、決められた処理をして出力を返す。 \\
\term{計算量}{Computational Complexity} & アルゴリズムが使う時間や記憶領域の増え方を表す尺度である。 \\
\term{基本ソフトウェア}{System Software} & ハードウェアを管理し、アプリケーションの実行基盤を提供するソフトウェアである。通常はオペレーティングシステムを指す。 \\
\term{プロトコル}{Protocol} & ネットワーク上で情報をやり取りするための約束事である。 \\
\term{人間要因}{Human Factors} & 利用者や開発者の認知、行動、作業しやすさがシステムに与える影響である。 \\
\bottomrule
\end{longtable}
